\clearpage\pagestyle{otMain}\setlength{\otBlocEspacement}{2.5mm}\fontsize{11}{15}\selectfont\corrigesChapitre{8}{Chercher des inconnues et des régularités}
\begin{blocLivre}{cor-8-1}\begin{corrige}{1}{Traduire par des barres}\demarche{Retirer l’écart, puis partager.}
On retire les 15 € d’écart : $75-15=60$ € pour deux barres égales. Lina a $\res{30}$ € et Marc $30+15=\res{45}$ €. Contrôle : $30+45=75$.
\label{fig-cor-8-1}\figureGeo[corrige]{alg-1}
\end{corrige}\end{blocLivre}
\begin{blocLivre}{cor-8-2}\begin{corrige}{2}{Quatre fois plus}\demarche{Compter les barres identiques.}
Cahier : 1 barre ; livre : 4 barres. Cinq barres valent 20 € : une barre vaut 4 €. Cahier : $\res4$ € ; livre : $\res{16}$ €.
\end{corrige}\end{blocLivre}
\begin{blocLivre}{cor-8-3}\begin{corrige}{3}{Poursuivre un motif}\demarche{Repérer ce qui s’ajoute à chaque étape.}
On passe de 4 à 7, puis à 10 : on ajoute 3 allumettes par carré. Étape 4 : $\res{13}$ ; étape 5 : $\res{16}$.
\label{fig-cor-8-3}\figureGeo[corrige]{alg-3}
\end{corrige}\end{blocLivre}
\begin{blocLivre}{cor-8-4}\begin{corrige}{4}{Le nombre pensé}\demarche{Défaire les étapes dans l’ordre inverse.}
On remonte : $39-7=32$, puis $32:4=\res8$. Contrôle : $8\times4+7=39$.
\end{corrige}\end{blocLivre}
\begin{blocLivre}{cor-8-5}\begin{corrige}{5}{Égalité de barres}\demarche{Enlever la partie connue, puis partager.}
Les trois barres valent $26-5=21$, donc une barre vaut $21:3=\res7$. Contrôle : $3\times7+5=26$.
\label{fig-cor-8-5}\figureGeo[corrige]{alg-5}
\end{corrige}\end{blocLivre}
\begin{blocLivre}{cor-8-6}\begin{corrige}{6}{Somme et écart}\demarche{Égaliser les barres en retirant l’écart.}
Deux barres égales valent $50-14=36$ : le petit nombre est $\res{18}$, le grand $\res{32}$. Contrôles : $18+32=50$ et $32-18=14$.
\end{corrige}\end{blocLivre}
\begin{blocLivre}{cor-8-7}\begin{corrige}{7}{Partage inégal}\demarche{Exprimer chaque part en barres de Bo.}
Bo : 1 barre ; Ali : 2 barres ; Cléa : 3 barres. Six barres valent 96 : une barre vaut 16. Bo $\res{16}$, Ali $\res{32}$, Cléa $\res{48}$.
\end{corrige}\end{blocLivre}
\begin{blocLivre}{cor-8-8}\begin{corrige}{8}{Le motif en L}\demarche{Un carreau d’angle, puis deux carreaux par étape.}
Étape 1 : 3 carreaux ; on ajoute 2 carreaux à chaque étape, un par branche. Étape 10 : $1+2\times10=\res{21}$. Pour 41 carreaux : $41-1=40$ et $40:2=$ étape $\res{20}$.
\label{fig-cor-8-8}\figureGeo[corrige]{alg-8}
\end{corrige}\end{blocLivre}
\begin{blocLivre}{cor-8-9}\begin{corrige}{9}{L’étape 100}\demarche{Décrire la structure plutôt que tout dessiner.}
Une allumette de départ, puis 3 par carré : $1+3\times100=\res{301}$ allumettes. Règle : « le nombre d’allumettes vaut 1 plus 3 fois le numéro de l’étape ».
\end{corrige}\end{blocLivre}
\begin{blocLivre}{cor-8-10}\begin{corrige}{10}{Les âges}\demarche{L’écart d’âge reste constant.}
Quatre barres valent 48 : la fille a $\res{12}$ ans et Paul $\res{36}$ ans. L’écart, 24 ans, ne change jamais. Paul a le double quand la fille a 24 ans : dans $\res{12}$ ans (Paul aura 48 ans).
\end{corrige}\end{blocLivre}
\begin{blocLivre}{cor-8-11}\begin{corrige}{11}{Les places de cinéma}\demarche{Ramener toutes les barres à la même taille.}
Si les adultes payaient le tarif enfant, on paierait $41-2\times3=35$ € pour 5 places : la place enfant coûte $\res7$ €, la place adulte $\res{10}$ €. Contrôle : $2\times10+3\times7=41$.
\end{corrige}\end{blocLivre}
\begin{blocLivre}{cor-8-12}\begin{corrige}{12}{L’escalier de cubes}\demarche{Repérer ce qui s’ajoute : de plus en plus de cubes.}
À l’étape $n$, on ajoute une colonne de $n$ cubes. Étape 4 : 10 ; étape 5 : 15 ; étape 6 : $\res{21}$. Étape 10 : $1+2+\ldots+10=\res{55}$.
\label{fig-cor-8-12}\figureGeo[corrige]{alg-12}
\end{corrige}\end{blocLivre}
\begin{blocLivre}{cor-8-13}\begin{corrige}{13}{La croix qui grandit}\demarche{Retirer la partie fixe, puis diviser.}
Un carreau central et 4 carreaux de plus à chaque étape : $1+4\times$ numéro. $101-1=100$ et $100:4=25$ : étape $\res{25}$.
\label{fig-cor-8-13}\figureGeo[corrige]{alg-13}
\end{corrige}\end{blocLivre}
\begin{blocLivre}{cor-8-14}\begin{corrige}{14}{Les deux réservoirs}\demarche{Le transfert réduit l’écart de deux fois 20 L.}
A : 3 barres ; B : 1 barre ; écart : 2 barres. Après le transfert, l’écart est nul : 2 barres valent $2\times20=40$ L. Une barre vaut 20 L : B contenait $\res{20}$ L et A $\res{60}$ L. Contrôle : $60-20=20+20$.
\end{corrige}\end{blocLivre}
\begin{blocLivre}{cor-8-15}\begin{corrige}{15}{Vrai ou faux}\demarche{Comparer avec la règle du motif.}
Étape 1 : 4 allumettes ; étape 10 : 31, et non $10\times4=40$. C’est $\res{\text{faux}}$ : l’allumette de départ n’est comptée qu’une fois.
\end{corrige}\end{blocLivre}
